\pdfbookmark[0]{Notices. Investigations into Topics in Higher Geodesy. 1843.}{anzeigen-geodaesie-erste}
\begin{center}
\rule{.85\textwidth}{.3pt}\par
Göttingische gelehrte Anzeigen.\quad 6 November 1843.\par
\rule{.85\textwidth}{.3pt}
\end{center}
\hspace*{\parindent}On 23 October, Court Councillor \textsc{Gauss} submitted a lecture to the Royal Society under the title
\begin{center}\textit{Investigations into Topics in Higher Geodesy},\end{center}
\noindent of which a short account will be given here.

In the trigonometrical part of the degree measurement that the author carried out in 1821---1827, and in the later trigonometrical surveys connected with it that extended over the whole Kingdom of Hanover, he took a different course from the usual one. This applied to the way the measurements were made, and even more to their subsequent mathematical treatment and the derivation of results from them. Some of Court Councillor \textsc{Gauss}'s distinctive methods have already been made public, partly by him in several papers published long ago and partly by others who learned of the methods in conversation or by letter and applied them in their own trigonometrical measurements. But the more substantial part of these methods has yet to be presented. These are the methods that differ most from those in common use and require a deeper mathematical grounding to be understood. The author had intended, once the measurements were entirely finished, to present the measurements themselves and all the procedures he had used in a separate work. Circumstances that need no explanation here have deferred that plan for the time being, so he has chosen instead to publish, in a series of separate memoirs, what is distinctive about his approach in the %
